%=============================================================================
% Section 11: Limitations and Evaluation Agenda
%=============================================================================
\section{Limitations and Evaluation Agenda}

This work focuses on architectural formalization and runtime boundaries. Beyond the evaluation agenda, we acknowledge several design tradeoffs and known limitations.

\subsection{Design Tradeoffs}

\textbf{Minimal vs.\ Full-Featured.} The AI Execution Layer deliberately constrains its responsibility to deterministic invocation. This minimalism enables portability and clear boundaries but sacrifices convenience features found in integrated frameworks:
\begin{itemize}
    \item \emph{No built-in policy logic.} Routing decisions, fallback chains, and cost optimization must be implemented by the Contact Layer, not the runtime itself.
    \item \emph{No conversation management.} Chat history, context windowing, and memory are outside the execution boundary.
    \item \emph{No built-in observability stack.} Metrics, tracing, and logging hooks are provided, but the observability infrastructure is deployment-specific.
\end{itemize}

\textbf{Protocol vs.\ Implementation.} The design prioritizes protocol stability over implementation convenience. Capability manifests must be declared explicitly; there is no ``auto-discovery by introspection'' mechanism. This tradeoff reduces runtime magic but increases provider onboarding effort.

\subsection{Known Boundaries}

The architecture is not suitable for all scenarios:
\begin{itemize}
    \item \emph{Single-provider lock-in.} Applications using provider-specific features beyond standardized capabilities lose portability benefits.
    \item \emph{Ultra-low-latency requirements.} The protocol translation layer adds overhead unsuitable for sub-millisecond response scenarios (e.g., high-frequency trading).
    \item \emph{Offline-only environments.} While edge runtimes are supported, capability discovery and provider registration assume network connectivity during initialization.
\end{itemize}

\subsection{Dependency Assumptions}

The architecture relies on several assumptions:
\begin{itemize}
    \item \emph{Protocol standardization.} The benefits diminish if providers do not adopt common capability schemas.
    \item \emph{Provider cooperation.} Capability manifests must accurately describe provider behavior; misaligned manifests cause runtime failures.
    \item \emph{Orchestration maturity.} Organizations without Contact Layer expertise may not fully realize the separation benefits.
\end{itemize}

\subsection{Evaluation Agenda}

A complete systems evaluation should include:
\begin{enumerate}
    \item \textbf{Translation Overhead:} protocol-to-provider transformation latency and variance.
    \item \textbf{Runtime Scalability:} throughput and tail latency under mixed synchronous/streaming workloads.
    \item \textbf{Failure Behavior:} bounded retry impact and error propagation under provider instability.
    \item \textbf{Policy-Execution Separation Cost:} control-plane/data-plane interaction overhead under dynamic routing.
\end{enumerate}

These measurements are left to subsequent experimental work and cross-runtime benchmark suites.
